\chapter{Eksperimentalni protokol}
\label{sec:eksperimenti}

V prejšnjih poglavjih smo definirali podatke, ciljne oznake, časovna okna,
grafovsko predstavitev in primerjane modele.
V tem poglavju določimo, kako te komponente ovrednotimo v skupni eksperimentalni
primerjavi.
Glavna evalvacijska naloga je binarna klasifikacija nizke proti visoki valenci.
Primerjava je zasnovana kot kombinacija modelov iz poglavja~\ref{sec:modeli} in
množic signalov iz razdelka~\ref{subsec:vhodni-signali}.
Vsi modeli uporabljajo iste 10-sekundne učne vzorce in iste delitve na učno,
validacijsko in testno množico.

\section{Zasnova glavne primerjave}
\label{subsec:zasnova-glavne-primerjave}

Cilj glavne primerjave je odgovoriti na raziskovalno vprašanje~\ref{rv:glavno}: \rvGlavnoBesedilo
Zato v primerjavi spreminjamo le način predstavitve okna in model, ki to predstavitev uporabi.

Ciljna oznaka je binarni razred valence, pri kateri model ločuje med nizko in visoko valenco.
Vse štiri množice iz tabele~\ref{tab:grafovska-konstrukcija-po-mnozicah} obravnavamo kot enakovreden del primerjave, saj nas zanima tudi, kako se klasifikacijska uspešnost spreminja z razpoložljivostjo različnih signalov sledilnika pogleda.

V primerjavi nastopajo vsi modeli iz tabele~\ref{tab:pregled-primerjanih-modelov}: izhodiščna klasifikatorja, klasični modeli nad agregiranimi značilkami, zamrznjeni prednaučeni predstavitveni modeli in grafovski modeli.
S tem primerjamo tri različne načine opisa istega časovnega okna: tabelarični povzetek signalov, vektorsko predstavitev časovne vrste in grafovsko predstavitev meritev.

\section{Validacijski protokol}
\label{subsec:validacijski-protokol}

\subsection{7-kratno prečno preverjanje po subjektih}
\label{subsec:subject-7fold}

Glavno primerjavo izvedemo s 7-kratnim prečnim preverjanjem po subjektih.
Pri vsaki ponovitvi so testni subjekti popolnoma ločeni od subjektov, ki so uporabljeni za učenje in validacijo.
Iz učnega dela izberemo eno skupino subjektov za validacijsko množico, preostali subjekti pa sestavljajo učno množico.
Delitve so pri vseh modelih in vseh množicah signalov enake, določene s fiksnim naključnim semenom.

Takšna delitev ocenjuje posploševanje na nove udeležence.
To je za uporabljeno nalogo pomembno, saj so signali sledilnika pogleda močno odvisni od specifik posameznika.
Ker 22 subjektov ni deljivo s sedem, imajo posamezne testne razdelitve različno velikost -- nekatere vsebujejo tri, druge pa štiri testne subjekte.

Validacijska množica se uporablja med učenjem, testna množica pa samo za končno oceno v posamezni ponovitvi.
Pri nevronskih modelih zato najboljšo različico modela izberemo glede na validacijsko izgubo, nato pa njeno delovanje izmerimo na pripadajoči testni množici.

\subsection{Predhodna preverjanja z LOO po subjektih}
\label{subsec:predhodna-loo-preverjanja}

V zgodnejših fazah smo za množico \emph{vsi signali} izvedli tudi LOO po subjektih.
Ker končna primerjava vključuje več modelov in štiri množice signalov, bi polni LOO protokol bistveno povečal število potrebnih eksperimentov.
Zato ga uporabimo kot podporno preverjanje, glavno primerjavo pa izvedemo s 7-kratnim protokolom.
Primerjava obeh validacijskih protokolov za valenco in množico \emph{vsi signali} je podana v dodatku~\ref{app:razsirjeni-rezultati-binarna-valenca}, kjer vidimo, da je 7-kratno prečno preverjanje primerljivo z metodo LOO.

\section{Nastavitve učenja in hiperparametri}
\label{subsec:nastavitve-ucenja}

V tem razdelku navedemo nastavitve učenja, ki vplivajo na izvedbo eksperimentov.
Nevronske modele učimo z zgodnjim ustavljanjem na validacijski izgubi.
Grafovski modeli uporabljajo operator \texttt{GCNConv}, dve grafovski plasti, skrite predstavitve velikosti 64 ter povezave s parametri $k_t=1$, $k_s=1$ in $k_f=3$, kot so definirani v poglavju~\ref{sec:grafovska-predstavitev}.
Podrobni hiperparametri vseh modelov so podani v dodatku~\ref{app:hiperparametri}.

Nastavitve grafovskih modelov izberemo na podlagi predhodnega kompaktnega mrežnega iskanja.
To iskanje je bilo izvedeno na 3-razredni valenci in ga v ~\ref{app:hiperparametri} poročamo samo kot pomožni poskus za izbiro stabilnih nastavitev.

\section{Metrike vrednotenja}
\label{subsec:metrike}

Glavni metriki vrednotenja sta točnost in makro F1.
Za vsak model in vsako množico signalov poročamo povprečje po sedmih testnih razdelitvah.
Točnost uporabimo kot osnovno metriko, ker ima glavna naloga dva enakovredna in globalno skoraj uravnotežena razreda, kot kaže slika~\ref{fig:porazdelitev-valenca-stevila}.

Makro F1 uporabimo kot dopolnilno glavno metriko.
Čeprav sta razreda globalno skoraj uravnotežena, so lahko posamezne testne razdelitve precej neuravnotežene (tudi do razmerja $1:2$).
Makro F1 izračuna F1 ločeno za vsak razred in nato razreda povpreči z enako utežjo.
S tem dopolni točnost predvsem pri modelih, ki lahko dosežejo sprejemljivo točnost, vendar večino primerov napovedujejo v en sam razred.

Poleg glavnih metrik spremljamo tudi uravnoteženo točnost, preciznost, priklic oziroma občutljivost, uteženi F1, AUC, izgubo in matrike zmede.
Te metrike uporabimo predvsem za preverjanje stabilnosti in za podrobnejšo analizo napak.
V glavnem besedilu se osredotočimo na točnost in makro F1, razširjene metrike pa poročamo v dodatku~\ref{app:razsirjeni-rezultati-binarna-valenca}. % če si podoben stavek napisal v poglavju 8, potem ga tu odstrani.

\section{Spremljanje učenja}
\label{subsec:spremljanje-ucenja}

Poleg končnih metrik shranjujemo tudi podatke o poteku učenja.
Pri nevronskih modelih spremljamo učno in validacijsko izgubo, najboljšo epoho, informacijo o zgodnjem ustavljanju, čas epohe in norme gradientov.
Krivulje izgub služijo predvsem preverjanju stabilnosti optimizacije in zaznavanju potencialnega prekomernega prileganja.

Pri grafovskih modelih spremljamo še nekaj diagnostičnih meritev naučenih predstavitev in napovedi.
Med njimi so varianca predstavitev grafa, povprečna kosinusna podobnost predstavitev, razpon logitov in entropija napovednih verjetnosti.
Te meritve uporabimo za preverjanje, ali se med učenjem pojavijo skoraj konstantne napovedi, zaznavanje kolapsa predstavitev ali znakov prekomernega glajenja.
Te diagnostike ne poročamo med rezultati v poglavju~\ref{sec:rezultati}, temveč le v dodatku~\ref{app:diagnostika-prekomerno-glajenje}.
